\chapter{Conclusiones y Líneas Futuras}

\section{Conclusiones}

El presente Trabajo de Fin de Grado tenía el objetivo de abordar el análisis de modelos de aprendizaje automático y aprendizaje profundo aplicados a la detección de intrusiones, considerando no solo su rendimiento predictivo, sino también su eficiencia computacional, su comportamiento frente a ataques adversarios y su posible viabilidad en escenarios de despliegue real. Podemos extraer, por tanto, varias conclusiones relevantes a partir de los resultados obtenidos.

En primer lugar, \textit{la optimización multiobjetivo mediante Optuna ha permitido seleccionar configuraciones de hiperparámetros atendiendo simultáneamente al F1-score y a la latencia de inferencia}. Esta parte ha sido especialmente relevante, ya que no se ha realizado una selección manual de modelos, sino un proceso sistemático de búsqueda para cada combinación de modelo y \textit{dataset}. El uso de la frontera de Pareto ha resultado útil para identificar configuraciones equilibradas, evitando seleccionar modelos excesivamente complejos que apenas aportan mejoras predictivas pero incrementan el coste computacional. De esta forma, la elección de los candidatos finales no se ha basado únicamente en maximizar el F1-score, sino en encontrar un equilibrio razonable entre capacidad de detección y eficiencia temporal.

En la comparativa entre modelos, \textit{los métodos basados en árboles han mostrado el comportamiento más equilibrado}. Random Forest, XGBoost, LightGBM y CatBoost han obtenido valores elevados de F1-score en los tres \textit{datasets}, manteniendo además una eficiencia temporal muy competitiva. Este resultado es especialmente importante porque confirma que, para datos tabulares de tráfico de red, los modelos de tipo \textit{ensemble} pueden ofrecer una combinación muy sólida entre precisión, latencia y \textit{throughput}. SVM ha destacado por su baja latencia y reducido coste de inferencia, aunque su rendimiento predictivo ha sido inferior al de los modelos no lineales, especialmente en UNSW-NB15. Por su parte, MLP y CNN-1D han alcanzado resultados competitivos, pero con un mayor coste de inferencia, especialmente en el caso de CNN-1D.

Los resultados también muestran diferencias importantes entre \textit{datasets}. En IoT-23 y TON-IoT los modelos han alcanzado valores de F1-score muy elevados en condiciones limpias, mientras que UNSW-NB15 ha presentado un escenario más exigente, con resultados más moderados. Esto confirma que \textit{el comportamiento de un IDS no depende únicamente del algoritmo utilizado, sino también del tipo de tráfico, la distribución de clases y las características propias de cada dataset}. Esta comparación entre conjuntos de datos ha permitido comprobar que un modelo puede comportarse de forma muy diferente según el escenario evaluado. Por tanto, evaluar únicamente sobre un dataset podría llevar a conclusiones incompletas o demasiado optimistas sobre la capacidad real del modelo.

El análisis frente a ataques de evasión mediante FGSM y PGD ha puesto de manifiesto que \textit{un buen rendimiento en test limpio no garantiza necesariamente robustez adversaria}. Algunos modelos que obtienen valores elevados de F1-score en condiciones normales sufren caídas importantes al evaluar muestras perturbadas. Además, el impacto de los ataques no es uniforme: depende del \textit{dataset}, del modelo atacado, del modelo utilizado para generar las perturbaciones y del límite de perturbación considerado. También se ha observado que el incremento de $\epsilon$ no siempre produce una degradación estrictamente monótona, especialmente en datos tabulares, donde perturbaciones elevadas pueden desplazar las muestras fuera de la distribución original o generar efectos de saturación. Esta parte del trabajo permite reforzar la idea de que la evaluación de un IDS no debe limitarse a datos limpios, ya que en un escenario real un atacante puede intentar modificar el tráfico para evadir la detección.

Desde el punto de vista computacional, los resultados de inferencia muestran que muchos de los modelos evaluados presentan un \textit{throughput} elevado, tanto en servidor como en un equipo local con menos recursos. Esto sugiere que, una vez entrenados y cargados en memoria, \textit{varios modelos clásicos podrían ser viables para escenarios de análisis en tiempo real}. Además, la comparación entre las mediciones realizadas en el servidor y en el equipo local ha permitido diferenciar entre el coste de un \textit{benchmark} instrumentado, donde se monitorizan recursos como CPU y RAM, y una medición centrada exclusivamente en inferencia pura.

En conjunto, los resultados obtenidos permiten concluir que \textit{la selección de un modelo IDS no debe basarse únicamente en maximizar una métrica predictiva}. Es necesario considerar de forma conjunta la distribución del \textit{dataset}, la capacidad de detección, la latencia, el \textit{throughput}, el consumo de recursos y la robustez frente a ataques adversarios.


\section{Líneas Futuras}

Esta línea de trabajo ha permitido obtener una visión integral del rendimiento y la eficiencia de distintos modelos de aprendizaje automático y profundo aplicados a la detección de intrusiones. Sin embargo, existen múltiples caminos de interés para continuar explorando y profundizando en esta área.

Incorporar y analizar modelos basados en secuencias temporales (como LSTM) podría ser especialmente relevante para capturar patrones de comportamiento a lo largo del tiempo, lo que es crucial en la detección de intrusiones que pueden manifestarse como anomalías temporales. 

Además, estudiar estrategias de reentrenamiento incremental permitiría evaluar cómo los modelos pueden adaptarse a nuevos datos sin necesidad de un reentrenamiento completo. Por ejemplo, se podría analizar con qué porcentaje de nuevos datos conviene reentrenar los modelos para mantener un buen rendimiento sin incurrir en un coste computacional excesivo.

Otro camino interesante sería aprovechar la similitud de características de los \textit{datasets} para evaluar la transferencia entre ellos, es decir, entrenar el modelo con un \textit{dataset} y probar la inferencia con otro. Esto permitiría analizar la capacidad de generalización de los modelos y su robustez.

Por último, una línea futura de gran interés sería analizar el entrenamiento para los ataques de evasión, reentrenando los modelos con muestras adversarias para mejorar su robustez frente a este tipo de amenazas, donde como se ha visto en el trabajo, los modelos pueden sufrir una degradación significativa de su eficacia.